\documentclass[10pt,a4paper,roman]{moderncv}      
\usepackage[english]{babel}
% \usepackage{hyperref}
% \hypersetup{
%     colorlinks=true,
%     linkcolor=blue,
%     filecolor=magenta,      
%     urlcolor=cyan,
%     pdftitle={Overleaf Example},
%     pdfpagemode=FullScreen,
%     }


\moderncvstyle{classic}                            
\moderncvcolor{green}                            

% character encoding
\usepackage[utf8]{inputenc}                     

% adjust the page margins
\usepackage[a4paper,
  total={6in, 8in},
  left=0.55in,
  right=0.55in,
  top=0.70in,
  bottom=0.25in,
  nohead
]{geometry}

% personal data
\name{Ming}{Wu}


\phone[mobile]{+32488390167}               
\email{ming.wu1@kuleuven.be}                             
\newcommand{\jobTitle}{\textbf{\textit{Assistant Professor in Mechanical Engineering
}}}

\newcommand{\jobInstitution}{\textit{University of Birmingham}}

\begin{document}
\recipient{To}{\jobInstitution
    }

\date{\today}
\opening{\textbf{Sub: Application for the \jobTitle}}
\closing{Your Sincerly, \vspace{-1em}}

\makelettertitle
Dear Hiring Committee,

I am writing to apply for the Assistant Professor position in Mechanical Engineering at the University of Birmingham. As a Research Associate at KU Leuven with a joint appointment in Manufacturing Processes & Systems (MaPS) and Declarative Languages & Artificial Intelligence (DTAI), I operate directly at the intersection of the two key areas identified in your search: advanced engineering and data science. My research establishes "Grounded Causal Intelligence"—a framework that integrates Causal AI with Adaptive Digital Twins to enable autonomous, zero-defect robotic manufacturing. I believe this expertise aligns perfectly with the School’s vision for "Phase 2" of its renaissance, particularly in the application of AI to Net-Zero and Future Mobility challenges.

\textbf{Research Vision:} AI for Net-Zero and Resilient Infrastructure The central hypothesis of my research is that the "physical residual"—the gap between idealized simulations and stochastic reality—contains the causal information necessary to eliminate trial-and-error in manufacturing. My work focuses on constructing Adaptive Digital Twins that actively query physical systems to correct physics-based models. This approach not only ensures "first-time-right" production for high-value components but also offers transferable frameworks for predictive maintenance in critical infrastructure.

I see a unique opportunity to scale this work within the University of Birmingham’s world-class ecosystem. Specifically, I aim to leverage the UK Rail Research and Innovation Network (UKRRIN) and the National Buried Infrastructure Facility to expand my "Adaptive Digital Twin" methodologies from component manufacturing to broader applications in transportation systems and structural health monitoring. This directly supports the department’s focus on sustainable transportation and data-centric engineering.

\textbf{Funding and Leadership} I bring a demonstrated track record of securing and managing competitive research funding. At KU Leuven, I serve as Lead Researcher for the €2.75M AutoCAM SBO project, where I direct the integration of generative AI with CAM planning. I also manage the technical execution of the €3.5M MUSIC SBO project, leading multi-sensor fusion for additive manufacturing. My current funding pipeline includes the AI4AM proposal (Horizon Europe, submitted Sept 2025), which aims to develop a Foundation Model for AM qualification, and the SYSTEMATIC project (Horizon Europe), focused on AI-driven Life Cycle Assessment (LCA) tools. I am committed to replicating this success at Birmingham by targeting EPSRC and Horizon Europe funding to support the School’s growth.

\textbf{Teaching:} Digital Engineering and the "Expert-in-the-Loop" I note the department’s specific requirement for teaching "digital engineering skills and AI methods." My teaching philosophy focuses on training engineers to be "Experts-in-the-Loop"—capable of using AI not just as a black box, but as a tool for hypothesis generation validated by physical first principles. I am prepared to deliver core undergraduate modules in Mechanics and Manufacturing Technology, alongside specialized courses in Applied AI for Engineering, Digital Twins, and Cyber-Physical Systems. My goal is to ensure Birmingham graduates possess the "bilingual" fluency in mechanics and data science required for Industry 4.0.

The University of Birmingham’s heritage of industrial collaboration and its bold investment in future-facing infrastructure make it the ideal environment for my research. I am eager to contribute to the department's leadership in sustainable manufacturing and intelligent transportation.

Thank you for your time and consideration. I would welcome the opportunity to discuss how my research can contribute to the School of Engineering.

\vspace{0.5cm}

\makeletterclosing

\end{document}
